\documentclass{article}
\usepackage{amsmath}
\usepackage{amssymb}
\usepackage{bm}
\usepackage{graphicx}
\usepackage{subcaption}
\usepackage{bigints}
\input{macros}
\begin{document}

\section*{Sheet G08 - Diffusion}
Robin Landsgesell, Jonas Marcic, Julian Stamm

\subsection*{Assignment 2}
\begin{align*}
	\gradient_Gf &= (df)^\sharp\\
	&= \left(\sum_{i=1}^n \frac{\partial f}{\partial x^i} dx^i\right)^\sharp\\
	&= \sum_{i=1}^n\sum_{j=1}^n g^{ij} \frac{\partial f}{\partial x^j} \frac{\partial}{\partial x^i}\\
\end{align*}

\subsection*{Assignment 3}

\subsubsection*{a)}
\begin{equation*}
	\frac{\partial S}{\partial r} = \begin{pmatrix}
		\cos(\varphi)\\
		\sin(\varphi)	
		\end{pmatrix},\quad
	\frac{\partial S}{\partial \varphi} = \begin{pmatrix}
		-r\sin(\varphi)\\
		r\cos(\varphi)	
		\end{pmatrix}
\end{equation*}\\
\begin{align*}
	g_{ij} &= \left\langle \frac{\partial S}{\partial x^i}, \frac{\partial S}{\partial x^j} \right\rangle\\
	g_{rr} &= \cos^2\varphi + \sin^2\varphi = 1\\
	g_{r\varphi} &= -r\cos\varphi\sin\varphi + r\sin\varphi\cos\varphi = 0\\
	g_{\varphi r} &= g_{r\varphi} = 0\\
	g_{\varphi \varphi} &= r^2\sin^2\varphi + r^2\cos^2\varphi = r^2\\
	\\
	G &= \begin{pmatrix}
		g_{rr} & g_{r\varphi}\\
		g_{\varphi r} & g_{\varphi \varphi}
		\end{pmatrix} = \begin{pmatrix}
		1 & 0\\
		0 & r^2
		\end{pmatrix}
\end{align*}

\newpage
\subsubsection*{b)}
\begin{center}
	$\boxed{\star dr}$\\
	\vspace{1em}
	Let $\star dr$ be of form $\star dr = a dr + b d\varphi$. For an arbitrary $\beta$ we get:
\end{center}

\begin{align*}
	\text{Compute LHS:}\\
	dr \wedge \beta &= \left\langle \star dr, \beta \right\rangle_\wedge dA\\
	dr \wedge (b_1dr+b_2d\varphi) &= \left\langle \star dr, \beta \right\rangle_\wedge dA\\
	b_2 dr \wedge d\varphi &= \left\langle \star dr, \beta \right\rangle_\wedge dA\\
	\text{Compute RHS:}\\
	b_2 dr \wedge d\varphi &= det(\left\langle \star dr, \beta \right\rangle_{G^{-1}}) dA\\
	b_2 dr \wedge d\varphi &= \begin{pmatrix}
								a\\
								b	
							  \end{pmatrix}^T\cdot G^{-1}\cdot \begin{pmatrix}
								b_1\\
								b_2	
								\end{pmatrix}dA\\
	b_2 dr \wedge d\varphi &= (ab_1 + \frac{1}{r^2}bb_2) \cdot dA\\
	\text{Plugin dA:}\\
	b_2 dr \wedge d\varphi &= (ab_1 + \frac{1}{r^2}bb_2) \cdot \sqrt{det G} \, dr \wedge d\varphi\\
	b_2 dr \wedge d\varphi &= (ab_1 + \frac{1}{r^2}bb_2) \cdot r dr \wedge d\varphi\\
	b_2 dr \wedge d\varphi &= (ab_1r + \frac{1}{r}bb_2)\, dr \wedge d\varphi\\
\end{align*}
\begin{center}
	We can conclude that:
\end{center}
\begin{equation*}
	\Rightarrow \begin{cases}
		ab_1r = 0\\
		\frac{1}{r}bb_2 = b_2
	\end{cases} \Rightarrow \begin{cases}
		a = 0\\
		b = r
	\end{cases} \Rightarrow \boxed{\star dr = r d\varphi}
\end{equation*}

\newpage
\begin{center}
	$\boxed{\star d\varphi}$\\
	\vspace{1em}
	Let $\star d\varphi$ be of form $\star d\varphi = a dr + b d\varphi$. For an arbitrary $\beta$ we get:
\end{center}

\begin{align*}
	\text{Compute LHS:}\\
	d\varphi \wedge \beta &= \left\langle \star d\varphi, \beta \right\rangle_\wedge dA\\
	d\varphi \wedge (b_1dr+b_2d\varphi) &= \left\langle \star d\varphi, \beta \right\rangle_\wedge dA\\
	b_1 d\varphi \wedge dr &= \left\langle \star d\varphi, \beta \right\rangle_\wedge dA\\
	\text{Compute RHS:}\\
	-b_1 dr \wedge d\varphi &= det(\left\langle \star d\varphi, \beta \right\rangle_{G^{-1}}) dA\\
	-b_1 dr \wedge d\varphi &= \begin{pmatrix}
								a\\
								b	
							  \end{pmatrix}^T\cdot G^{-1}\cdot \begin{pmatrix}
								b_1\\
								b_2	
								\end{pmatrix}dA\\
	-b_1 dr \wedge d\varphi &= (ab_1 + \frac{1}{r^2}bb_2) \cdot dA\\
	\text{Plugin dA:}\\
	-b_1 dr \wedge d\varphi &= (ab_1 + \frac{1}{r^2}bb_2) \cdot \sqrt{det G} \, dr \wedge d\varphi\\
	-b_1 dr \wedge d\varphi &= (ab_1 + \frac{1}{r^2}bb_2) \cdot r dr \wedge d\varphi\\
	-b_1 dr \wedge d\varphi &= (ab_1r + \frac{1}{r}bb_2)\, dr \wedge d\varphi\\
\end{align*}
\begin{center}
	We can conclude that:
\end{center}
\begin{equation*}
	\Rightarrow \begin{cases}
		ab_1r = -b_1\\
		\frac{1}{r}bb_2 = 0
	\end{cases} \Rightarrow \begin{cases}
		a = -\frac{1}{r}\\
		b = 0
	\end{cases} \Rightarrow \boxed{\star d\varphi = -\frac{1}{r} dr}
\end{equation*}

\newpage
\begin{center}
	$\boxed{\star (dr \wedge d\varphi)}$\\
	\vspace{1em}
	Let $\star (dr \wedge d\varphi)$ be of form $\star (dr \wedge d\varphi) = c$. Because $(dr \wedge d\varphi)$ is a 2-form, so $\star(dr \wedge d\varphi)$ has to be a 0-form. Let $\beta = (b_1dr + b_2 d\varphi + b_3)$ and
	$G^{-1} = \begin{pmatrix}
		1 & 0 & 0\\
		0 & \frac{1}{r^2} & 0\\
		0 & 0 & 1
		\end{pmatrix}$:
\end{center}
\begin{align*}
	(dr \wedge d\varphi) \wedge \beta &= \left\langle \star (dr \wedge d\varphi), \beta \right\rangle_\wedge dA\\
	(dr \wedge d\varphi) \wedge (b_1dr+b_2d\varphi + b_3) &= \left\langle \star (dr \wedge d\varphi), \beta \right\rangle_\wedge dA\\
	b_3\, dr \wedge d\varphi &= \left\langle c, \beta \right\rangle_\wedge dA\\
	b_3\, dr \wedge d\varphi &= det(\left\langle c, \beta \right\rangle_{G^{-1}}) dA\\
	b_3\, dr \wedge d\varphi &= \begin{pmatrix}
								0\\
								0\\
								c	
							  \end{pmatrix}^T\cdot G^{-1}\cdot \begin{pmatrix}
								b_1\\
								b_2\\
								b_3
								\end{pmatrix}dA\\
	b_3\, dr \wedge d\varphi &= c \cdot b_3 \cdot dA\\
	b_3\, dr \wedge d\varphi &= c \cdot b_3 \sqrt{det G} \, dr \wedge d\varphi\\
	b_3\, dr \wedge d\varphi &= c \cdot b_3 r\, dr \wedge d\varphi\\
\end{align*}
\begin{center}
	We can conclude that:
\end{center}
\begin{equation*}
	\Rightarrow c \cdot b_3 r = b_3 \Rightarrow c = \frac{1}{r} \Rightarrow \boxed{\star (dr \wedge d\varphi) = \frac{1}{r}}
\end{equation*}

\subsubsection*{c)}
1. Step: Compute $df$:
\begin{equation*}
	df = \frac{\partial f}{\partial r} dr + \frac{\partial f}{\partial \varphi} d\varphi
\end{equation*}
2. Step: Compute $\star df$:
\begin{align*}
	\star df &= \star \left(\frac{\partial f}{\partial r} dr + \frac{\partial f}{\partial \varphi} d\varphi\right)\\
	&= \frac{\partial f}{\partial r} \star dr + \frac{\partial f}{\partial \varphi} \star d\varphi\\
	&= \frac{\partial f}{\partial r} (r d\varphi) + \frac{\partial f}{\partial \varphi} \left(-\frac{1}{r} dr\right)\\
	&= r \frac{\partial f}{\partial r} d\varphi - \frac{1}{r} \frac{\partial f}{\partial \varphi} dr
\end{align*}
3. Step: Compute $d \star df$:
\begin{align*}
	d \star df &= d\left(r \frac{\partial f}{\partial r} d\varphi - \frac{1}{r} \frac{\partial f}{\partial \varphi} dr\right)\\
	&= d\left(r \frac{\partial f}{\partial r} d\varphi\right) - d\left(\frac{1}{r} \frac{\partial f}{\partial \varphi} dr\right)\\
	&= \frac{\partial}{\partial r}\left(r \frac{\partial f}{\partial r}\right) dr \wedge d\varphi + \frac{\partial}{\partial \varphi}\left(- \frac{1}{r} \frac{\partial f}{\partial \varphi}\right) d\varphi \wedge dr\\
	&= \left[\frac{\partial}{\partial r}\left(r \frac{\partial f}{\partial r}\right) + \frac{\partial}{\partial \varphi}\left(\frac{1}{r} \frac{\partial f}{\partial \varphi}\right)\right] dr \wedge d\varphi
\end{align*}
4. Step: Compute $\star d \star df$:
\begin{align*}
	\star d \star df &= \star \left[\left(\frac{\partial}{\partial r}\left(r \frac{\partial f}{\partial r}\right) + \frac{\partial}{\partial \varphi}\left(\frac{1}{r} \frac{\partial f}{\partial \varphi}\right)\right) dr \wedge d\varphi\right]\\
	&= \left(\frac{\partial}{\partial r}\left(r \frac{\partial f}{\partial r}\right) + \frac{\partial}{\partial \varphi}\left(\frac{1}{r} \frac{\partial f}{\partial \varphi}\right)\right) \star (dr \wedge d\varphi)\\
	&= \left(\frac{\partial}{\partial r}\left(r \frac{\partial f}{\partial r}\right) + \frac{\partial}{\partial \varphi}\left(\frac{1}{r} \frac{\partial f}{\partial \varphi}\right)\right) \cdot \frac{1}{r}\\
	&= \frac{1}{r} \frac{\partial}{\partial r}\left(r \frac{\partial f}{\partial r}\right) + \frac{1}{r} \frac{\partial}{\partial \varphi}\left( \frac{1}{r}\frac{\partial f}{\partial \varphi}\right)\\
	&= \frac{1}{r} \frac{\partial}{\partial r}\left(r \frac{\partial f}{\partial r}\right) + \frac{1}{r^2} \left( \frac{\partial^2 f}{\partial \varphi^2}\right)\\
	&= \Delta f
\end{align*}

\end{document}